Graph theoretic activities allow us to perform tasks on graphs.
Tables~\ref{tab:graphactivities:1}-\ref{tab:graphactivities:2}
show the commands for graph theoretic activities.
These are activities that can be applied to objects of type graph. 
\orbitertableofcommands{graph_theoretic_activities_1}
\orbitertableofcommands{graph_theoretic_activities_2}


\bigskip

Let us continue the example of the three-cycle 
from Section~\ref{sec:graphtheory}. 
The command 

\sourcecode{triangle_graph_properties}

computes the degree type of the graph. 
This is the distribution of degrees in the degree sequence of the graph.
It prints the distribution of degree values in exponential notation. 
The multiplicities are indicated as exponent. 
For instance, the graph in this example has three vertices of degree 2, 
so the degree type is written as $2^3.$

\bigskip


The command 

\sourcecode{Hamming_graph_7_drg}

verifies that the Hamming graph of order $7$ is distance regular. 
The command assumes that the graph is known to be vertex transitive. 
The computation is performed using vertex $0$ as a distinguished vertex. 
The vertex for which the computation is to be performed is an argument to the command.
If the graph is distance regular, the parameters $a_i,b_i,c_i$ are computed 
for $i =0,\ldots,d$, where $d$ is the diameter of the graph.


\bigskip

The command 

\sourcecode{Johnson_8_3_2_drg}

verifies that the Johnson graph $J_{8,3}$ is distance regular. 



\bigskip


The following command
exports the adjacency matrix and creates a bitmap drawing of it.

\sourcecode{Cycle_13_draw}

The command first creates the cycle graph of order 13, 
and then exports the adjacency matrix as 
csv file. 
It then draws the adjacency matrix as a bitmap graphics. 

\bigskip

Suppose we want to compute the eigenvalues 
of the adjacency matrix of a graph.
In the following example, 
the command \verb'-eigenvalues' is used to compute 
the eigenvalues (both regular and Laplace) 
of the 9-cycle: 

\sourcecode{Cycle_9_eigenvalues}

The following output is produced:
\orbiteroutput{GRAPHICS/LATEX/eigenvalues_Cycle_9}


\bigskip



The command 

\sourcecode{Paley_13_draw}

draws the Paley graph of order 13 created 
in Section~\ref{sec:graphtheory} 
using the external tool graphviz.

\bigskip

Let us consider the Cayley graphs from Section~\ref{sec:graphtheory}.
Here is a command that draws the first graph and computes the eigenvalues:

\sourcecode{Cayley_Z11_1mod3_eigenvalues_and_draw}

The drawing is shown below
\orbiteroutput{GRAPHICS/WRAPPERS/graph_Cayley_Z11_1_and_3}
Let us draw the Cayley graph of $\Sym(5)$ 
with respect to the Coxeter generators:

\sourcecode{Cayley_Sym5_coxeter_draw}

The drawing is shown below
\orbiteroutput{GRAPHICS/WRAPPERS/graph_Cayley_Sym5_Coxeter}



\bigskip


It is possible to create the collinearity graph of an incidence structure. 
The incidence structure must be encoded by means of an incidence matrix. 
We can search for disjoint sets in the generalized quadrangle $W(2)$ (the Doily).
For starters, here is the makefile variable to create the geometry:

\sourcecodevariable{DOILY}

Next, we read the incidence matrix from file and create the collinearity graph of it:


\sourcecode{PGO_5_2_collinearity_graph}

The command also computes properties of the graph. 
The graph has 15 vertices of degree 6. 
This is because in the geometry, each point lies on three lines, 
and hence is collinar with 6 other points.




\bigskip

The command 

\sourcecode{Op_4_2_affine_polar_graph_info}

computes properties of the affine polar graph based on $O^+(4,2).$

\bigskip

The command 

\sourcecode{Neumaier_16_test_SRG_property}

tests the Neumaier property for the Neumaier graph on 16 vertices.  

\bigskip


The command 

\sourcecode{Paulus_1_aut}

computes the automorphism group of 
the first graph in the list of Paulus and Rozenfeldt.


\bigskip

The command 

\sourcecode{Paulus_15_aut}

computes the automorphism group of 
the 15th graph in the list of Paulus and Rozenfeldt.


\bigskip



The command 

\sourcecode{PG_3_2_packing_make_group_and_test}

creates the set of all labeled spreads in $\PG(3,2).$ 
It then creates the graph of disjointness on the set of spreads. 
Afterwards, it creates the action of the projective group on the set of all labeled spreads in $\PG(3,2).$ 
Finally, it tests whether the group preserves the graph. 


\bigskip




The command 

\sourcecode{Petersen_graph_distance_poset}

computes the poset of neighborhood sets of the Petersen graph with respect to vertex 0.

The command 

\sourcecode{Petersen_graph_distance_poset_draw}

produces a drawing of the poset, shown below:

\orbiteroutput{GRAPHICS/WRAPPERS/Petersen_graph_distance_poset_draw}

The command 

\sourcecode{Petersen_test_if_distance_regular}

tests whether the Petersen graph is distance regular. The answer is positive.

\bigskip


The command 

\sourcecode{Cayley_Sym5_distance_poset}

computes the neighborhood poset for the Coxeter group $\Sym(5).$ 
This poset is related to the Bruhat order.
The following drawing is produced:

\orbiteroutput{GRAPHICS/WRAPPERS/Cayley_Sym5_distance_poset}






